\documentclass[11pt]{article}
\usepackage[margin=1.1in]{geometry}
\usepackage{amsmath,amssymb,amsthm,mathtools,booktabs,array,tabularx,longtable}
\usepackage[dvipsnames]{xcolor}
\usepackage[colorlinks=true,linkcolor=MidnightBlue,citecolor=MidnightBlue,urlcolor=MidnightBlue]{hyperref}
\usepackage{fancyhdr}
\pagestyle{fancy}\fancyhf{}\renewcommand{\headrulewidth}{0pt}
\fancyfoot[R]{\thepage}
\fancyfoot[L]{\scriptsize\color{gray}Generated by the Transport Lab from approved parts of its store; not a deposited paper. Papers of record: \href{https://zenodo.org/communities/transport-papers}{Zenodo, Transport Papers}}
\fancypagestyle{plain}{\fancyhf{}\fancyfoot[R]{\thepage}\fancyfoot[L]{\scriptsize\color{gray}Generated by the Transport Lab from approved parts of its store; not a deposited paper. Papers of record: \href{https://zenodo.org/communities/transport-papers}{Zenodo, Transport Papers}}\renewcommand{\headrulewidth}{0pt}}
\newtheorem{theorem}{Theorem}[section]
\newtheorem{proposition}[theorem]{Proposition}
\newtheorem{lemma}[theorem]{Lemma}
\newtheorem{corollary}[theorem]{Corollary}
\newtheorem{conjecture}[theorem]{Conjecture}
\newtheorem{refutation}[theorem]{Refutation}
\newtheorem{prediction}[theorem]{Prediction}
\newtheorem{principle}[theorem]{Principle}
\newtheorem{observation}[theorem]{Observation}
\newtheorem{axiom}[theorem]{Axiom}
\theoremstyle{definition}
\newtheorem{definition}[theorem]{Definition}
\newtheorem{fact}[theorem]{Fact}
\newtheorem{assembled}[theorem]{Assembly}
\newtheorem{postulate}[theorem]{Postulate}
\newtheorem{premise}[theorem]{Premise}
\newtheorem{hypothesis}[theorem]{Hypothesis}
\newtheorem{remark}[theorem]{Remark}
\newtheorem{example}[theorem]{Example}
\theoremstyle{plain}
\setlength{\marginparwidth}{0.8in}
\newcommand{\changedmark}[2]{\marginpar{\raggedright\scriptsize\color{#1}#2}}
\providecommand{\Aa}{}\renewcommand{\Aa}{\mathcal{A}}
\providecommand{\Ad}{}\renewcommand{\Ad}{\widetilde{\mathrm{Ad}}}
\providecommand{\Aev}{}\renewcommand{\Aev}{\mathcal{A}^{+}}
\providecommand{\Aod}{}\renewcommand{\Aod}{\mathcal{A}^{-}}
\providecommand{\B}{}\renewcommand{\B}{\mathrm{B}}
\providecommand{\C}{}\renewcommand{\C}{\mathbb{C}}
\providecommand{\CAR}{}\renewcommand{\CAR}{\mathrm{CAR}}
\providecommand{\CS}{}\renewcommand{\CS}{{\mathbb{S}^{+}}}
\providecommand{\CSm}{}\renewcommand{\CSm}{{\mathbb{S}^{-}}}
\providecommand{\Cl}{}\renewcommand{\Cl}{\mathrm{Cl}}
\providecommand{\Con}{}\renewcommand{\Con}{\mathcal{C}}
\providecommand{\Dl}{}\renewcommand{\Dl}{\Delta}
\providecommand{\EN}{}\renewcommand{\EN}{E_{N}}
\providecommand{\End}{}\renewcommand{\End}{\mathrm{End}}
\providecommand{\F}{}\renewcommand{\F}{\mathbb{F}}
\providecommand{\G}{}\renewcommand{\G}{\Gamma}
\providecommand{\HH}{}\renewcommand{\HH}{\mathbb{H}}
\providecommand{\IS}{}\renewcommand{\IS}{I\sigma_3}
\providecommand{\J}{}\renewcommand{\J}{\mathbf{J}}
\providecommand{\LH}{}\renewcommand{\LH}{L_{H}}
\providecommand{\Mat}{}\renewcommand{\Mat}{\mathrm{M}}
\providecommand{\Mo}{}\renewcommand{\Mo}{M^{\circ}}
\providecommand{\Mp}{}\renewcommand{\Mp}{M^{+}}
\providecommand{\Msun}{}\renewcommand{\Msun}{M_{\odot}}
\providecommand{\Mt}{}\renewcommand{\Mt}{\widetilde{M}}
\providecommand{\Ncone}{}\renewcommand{\Ncone}{\mathcal{N}}
\providecommand{\Neff}{}\renewcommand{\Neff}{N_{\mathrm{eff}}}
\providecommand{\PCS}{}\renewcommand{\PCS}{\mathbb{P}{\mathbb{S}^{+}}}
\providecommand{\PT}{}\renewcommand{\PT}{\mathbb{PT}}
\providecommand{\Pin}{}\renewcommand{\Pin}{\mathrm{Pin}}
\providecommand{\Q}{}\renewcommand{\Q}{\mathbb{Q}}
\providecommand{\R}{}\renewcommand{\R}{\mathbb{R}}
\providecommand{\Ric}{}\renewcommand{\Ric}{\mathrm{Ric}}
\providecommand{\SBH}{}\renewcommand{\SBH}{S_{\mathrm{BH}}}
\providecommand{\Sc}{}\renewcommand{\Sc}{\mathbb{S}}
\providecommand{\Sch}{}\renewcommand{\Sch}{\mathrm{P}}
\providecommand{\Sig}{}\renewcommand{\Sig}{\Sigma}
\providecommand{\Som}{}\renewcommand{\Som}{S_\omega}
\providecommand{\Sp}{}\renewcommand{\Sp}{\mathbb{S}}
\providecommand{\Sph}{}\renewcommand{\Sph}{\mathrm{S}}
\providecommand{\Spi}{}\renewcommand{\Spi}{S_\pi}
\providecommand{\Spin}{}\renewcommand{\Spin}{\mathrm{Spin}}
\providecommand{\Sth}{}\renewcommand{\Sth}{S_{\mathrm{th}}}
\providecommand{\T}{}\renewcommand{\T}{\mathbb{T}}
\providecommand{\TH}{}\renewcommand{\TH}{T_{\mathrm{H}}}
\providecommand{\Tr}{}\renewcommand{\Tr}{\operatorname{Tr}}
\providecommand{\UHF}{}\renewcommand{\UHF}{\mathrm{UHF}}
\providecommand{\Z}{}\renewcommand{\Z}{\mathbb{Z}}
\providecommand{\Zc}{}\renewcommand{\Zc}{\mathrm{Z}}
\providecommand{\cmg}{}\renewcommand{\cmg}{\mathrm{cm^{2}\,g^{-1}}}
\providecommand{\dS}{}\renewcommand{\dS}{\mathrm{dS}}
\providecommand{\dsc}{}\renewcommand{\dsc}{\mathrm{dsc}}
\providecommand{\dtPG}{}\renewcommand{\dtPG}{dt}
\providecommand{\dtS}{}\renewcommand{\dtS}{dt_{\mathrm{S}}}
\providecommand{\ellv}{}\renewcommand{\ellv}{\ell}
\providecommand{\er}{}\renewcommand{\er}{e_r}
\providecommand{\et}{}\renewcommand{\et}{e_t}
\providecommand{\ev}{}\renewcommand{\ev}{\mathrm{ev}}
\providecommand{\gl}{}\renewcommand{\gl}{\mathfrak{gl}}
\providecommand{\hatV}{}\renewcommand{\hatV}{\widehat{V}}
\providecommand{\hb}{}\renewcommand{\hb}{\bar{h}}
\providecommand{\hf}{}\renewcommand{\hf}{\hat{f}}
\providecommand{\id}{}\renewcommand{\id}{\mathrm{id}}
\providecommand{\kpc}{}\renewcommand{\kpc}{\mathrm{kpc}}
\providecommand{\lP}{}\renewcommand{\lP}{\ell_{P}}
\providecommand{\mD}{}\renewcommand{\mD}{\mathbf{m}_{D}}
\providecommand{\mb}{}\renewcommand{\mb}{\bar m}
\providecommand{\mnu}{}\renewcommand{\mnu}{\Sigma m_{\nu}}
\providecommand{\nInf}{}\renewcommand{\nInf}{n_\infty}
\providecommand{\nO}{}\renewcommand{\nO}{n_0}
\providecommand{\ninf}{}\renewcommand{\ninf}{n_{\infty}}
\providecommand{\nz}{}\renewcommand{\nz}{n_{0}}
\providecommand{\proj}{}\renewcommand{\proj}{\operatorname{proj}}
\providecommand{\pv}{}\renewcommand{\pv}{\pi}
\providecommand{\rad}{}\renewcommand{\rad}{\mathrm{rad}}
\providecommand{\sgn}{}\renewcommand{\sgn}{\operatorname{sgn}}
\providecommand{\shA}{}\renewcommand{\shA}{\mathcal{A}}
\providecommand{\shC}{}\renewcommand{\shC}{\mathcal{C}}
\providecommand{\shE}{}\renewcommand{\shE}{\mathcal{E}}
\providecommand{\shJ}{}\renewcommand{\shJ}{\mathcal{J}}
\providecommand{\shK}{}\renewcommand{\shK}{\mathcal{K}}
\providecommand{\shL}{}\renewcommand{\shL}{L}
\providecommand{\shW}{}\renewcommand{\shW}{\mathcal{W}}
\providecommand{\sll}{}\renewcommand{\sll}{\mathfrak{sl}}
\providecommand{\so}{}\renewcommand{\so}{\mathfrak{so}}
\providecommand{\su}{}\renewcommand{\su}{\mathfrak{su}}
\providecommand{\supp}{}\renewcommand{\supp}{\operatorname{supp}}
\providecommand{\tr}{}\renewcommand{\tr}{\operatorname{tr}}
\providecommand{\uu}{}\renewcommand{\uu}{\mathfrak{u}}
\providecommand{\br}{}\renewcommand{\br}[1]{\langle #1\rangle}
\providecommand{\bra}{}\renewcommand{\bra}[1]{\langle #1|}
\providecommand{\file}{}\renewcommand{\file}[1]{\texttt{#1}}
\providecommand{\hp}{}\renewcommand{\hp}[2]{( #1 \,|\, #2 )}
\providecommand{\ip}{}\renewcommand{\ip}[2]{\langle #1, #2\rangle}
\providecommand{\ket}{}\renewcommand{\ket}[1]{|#1\rangle}
\providecommand{\rel}{}\renewcommand{\rel}[2]{D\bigl(#1 \,\|\, #2\bigr)}
\providecommand{\rev}{}\renewcommand{\rev}[1]{\widetilde{#1}}
\providecommand{\slashed}{}\renewcommand{\slashed}[1]{{#1}\mkern-10mu/}
\title{\bfseries The completion}
\author{Leon Butler\\ \normalsize\href{mailto:lbutler.math@gmail.com}{\texttt{lbutler.math@gmail.com}}\\ \normalsize ORCID: \href{https://orcid.org/0009-0006-4449-5617}{0009-0006-4449-5617}}
\date{Generated by the Transport Lab, 2 November 2026}
\begin{document}
\maketitle
\begin{abstract}
\normalsize
When division fails, the series does not give up: it moves to a larger setting in which the failing calculation works. This trail follows that move, the completion.

For a null vector the larger space must gain one new positive and one new negative direction, so the minimal completion is $\Cl(p,q) \to \Cl(p{+}1, q{+}1)$. There a single operation gives a reciprocal everywhere, agreeing with the usual one wherever that exists.

For an algebra split into two halves, adding one direction of either sign fuses the halves. Together these give the transport principle: three kinds of failure, each with a forced smallest repair.
\end{abstract}
\newpage
\noindent \textbf{Authorship and disclosure. }The founding idea of the series is the author's: that division by zero in a Clifford algebra should be read not as a prohibition but as a transport into a higher-dimensional algebra, and that the transported division need not be accompanied by an inverse operation. The direction of investigation and the critical review of the results are likewise the author's. The mathematical development and the drafting of the text were carried out in extended collaboration with Claude Fable 5 (Anthropic), an AI system.

\noindent\fbox{\parbox{\dimexpr\linewidth-2\fboxsep-2\fboxrule}{\textbf{This is not a deposited paper.} It was generated by the Transport Lab on 2 November 2026 and assembled from approved parts of its store. It follows results across several papers of the series; each is stated with its proof and credited to the paper that first proved it. The official versions of those papers are on Zenodo, in the Transport Papers community, \url{https://zenodo.org/communities/transport-papers}.}}
\newpage

\section*{Introduction}
The walk first repairs null vectors on the space of points, then degenerate geometries, then algebras split in two, and ends with the principle that gathers the repairs together.

\section{The completion}

Each result here says how to enlarge a space or an algebra so that a failure of division is repaired, and how small that enlargement can be.

A null vector cannot be divided by, so the series instead represents each point of the geometry as a null vector in a larger space. The first question is how much larger that space must be.

The added part needs at least one direction of positive square and one of negative square, and one such pair, a hyperbolic plane, is enough: \[ \Cl(p,q) \;\longrightarrow\; \Cl(p{+}1, q{+}1) \] with signature increment $(+1,+1)$ in every signature.

\begin{theorem}[Minimality]\label{c:transport.minimality}
To make the reciprocal work everywhere, the series represents each point as a light-like direction in a larger space. This result asks how much larger that space has to be. The added part must contain at least one direction of positive square and one of negative square, so adding only space-like or only time-like directions never works, while one such pair is enough and is essentially unique.
\emph{First proved in} T1, Thm 4.13 (\href{https://doi.org/10.5281/zenodo.21743863}{doi:10.5281/zenodo.21743863}).
\end{theorem}
\noindent\emph{Rests on:} \ref{c:std-algebra.witt-decomposition}, \ref{c:std-algebra.sylvester-inertia}.

With the larger space in hand, the reciprocal can be extended.

On the space of points a single sandwich by the new direction $e_+$, \[ \iota[N] = [\,e\_+ N e\_+\,] \] gives a reciprocal everywhere: on vectors with $Q(x) \neq 0$ it is the usual $x/Q(x)$, and it sends null vectors to points at infinity. It is also the only continuous rule that matches the usual reciprocal.

\begin{theorem}[Totality]\label{c:transport.totality}
In a Clifford algebra with vectors of zero length, the reciprocal of a vector is undefined on exactly those vectors. After moving to the enlarged space of points, a single reflection-like operation gives a reciprocal everywhere: it agrees with the usual one on ordinary vectors and sends zero-length vectors to points at infinity. Those values are the limits one would expect, and no other continuous rule matches the usual reciprocal.
\emph{First proved in} T1, Thm 4.10 (\href{https://doi.org/10.5281/zenodo.21743863}{doi:10.5281/zenodo.21743863}).
\end{theorem}
\noindent\emph{Rests on:} \ref{c:transport.density}, \ref{c:std-algebra.reflection}.

A calculation done in the larger space is only useful if its answer can be brought back.

The read-out map \[ \dsc\,[\, s\,\nO + x + t\,\nInf \,] = x/s \] does this for every point not at infinity, and on invertible vectors it returns the familiar reciprocal $x/Q(x)$. For null vectors the answer lies at infinity, and the read-out declines rather than inventing a value.

\begin{lemma}[Descent: the read-out map]\label{c:transport.descent-read-out-map}
After computing in the enlarged space one needs a way back to the original one. The read-out map does this for every point not at infinity, and for ordinary invertible vectors it gives back the familiar reciprocal. For zero-length vectors the answer lies at infinity, and the read-out declines, because no point of the original space represents it.
\emph{First proved in} T1, Lem 4.23 (\href{https://doi.org/10.5281/zenodo.21743863}{doi:10.5281/zenodo.21743863}).
\end{lemma}
\noindent\emph{Rests on:} \ref{c:transport.totality}.

The same idea repairs geometries whose form is degenerate, with directions that pair to zero with everything.

Each degenerate direction needs exactly one new partner, so a radical of dimension $r$ gives \[ \Cl(p,q,r) \hookrightarrow \Cl(p{+}r, q{+}r) \] unique up to an isometry fixing the original space. Minimality here is among enlargements of the underlying space, not among all algebras.

\begin{theorem}[Minimal nondegenerate completion]\label{c:conformal.minimal-nondegenerate-completion}
A geometric algebra built on a degenerate metric has directions that pair to zero with everything, and division fails along them. This result says how to enlarge the space so that the metric becomes nondegenerate: each degenerate direction needs exactly one new partner direction, and the enlarged metric gains one positive and one negative direction for each degenerate one. The smallest such enlargement is unique up to a symmetry that leaves the original space untouched. Minimality here is among enlargements of the underlying space, not among all algebras.
\emph{First proved in} T2, Thm 2.3 (\href{https://doi.org/10.5281/zenodo.22714886}{doi:10.5281/zenodo.22714886}).
\end{theorem}
\noindent\emph{Rests on:} \ref{c:conformal.blade-basis-persists-under-degeneracy}, \ref{c:conformal.witt-extension-theorem}, \ref{c:std-algebra.functoriality}.

The best-known example comes from engineering and computer graphics.

Projective geometric algebra, with one degenerate direction for the plane at infinity, completes uniquely to conformal geometric algebra: \[ \Cl(3,0,1) \;\longrightarrow\; \Cl(4,1) \] and in general $\Cl(n,0,1)$ completes to $\Cl(n{+}1,1)$. Two tools usually treated as separate design choices are thus one algebra and its repair.

\begin{corollary}[The principal case]\label{c:conformal.principal-case}
Engineers and graphics programmers use two geometric algebras: projective geometric algebra, which has one degenerate direction standing for the plane at infinity, and conformal geometric algebra, which has two extra directions. They are usually treated as separate design choices. This result shows that the conformal algebra is the unique smallest nondegenerate enlargement of the projective one, in any number of dimensions, in the sense of enlarging the underlying space.
\emph{First proved in} T2, Cor 3.1 (\href{https://doi.org/10.5281/zenodo.22714886}{doi:10.5281/zenodo.22714886}).
\end{corollary}
\noindent\emph{Rests on:} \ref{c:conformal.minimal-nondegenerate-completion}.

Back with the null repair, it helps to know what the enlarged algebra actually is.

Adding one hyperbolic plane, the step called Rule 1, gives two-by-two matrices over the original algebra, \[ \Cl(p{+}1, q{+}1) \;\cong\; \Mat\_2\bigl(\Cl(p,q)\bigr) \] with the new null pair $f, f^*$ as the off-diagonal entries. So the step leaves the clock residue $d = p - q \pmod 8$ unchanged.

\begin{theorem}[Rule 1 is Morita-trivial]\label{c:transport.rule-1-morita-trivial}
The basic transport step adds one new pair of null directions to an algebra. The enlarged algebra turns out to be the algebra of two-by-two matrices with entries in the original one, with the new null pair playing the role of the off-diagonal entries. So the step leaves the essential type of the algebra unchanged, although whether certain other kinds of obstruction are present can change.
\emph{First proved in} T1, Thm 7.1 (\href{https://doi.org/10.5281/zenodo.21743863}{doi:10.5281/zenodo.21743863}).
\end{theorem}
\noindent\emph{Rests on:} \ref{c:transport.matrix-units-peirce-isomorphism}, \ref{c:std-algebra.involution-properties}, \ref{c:std-algebra.morita-amplification}.

The second kind of failure, an algebra split into two halves, needs a different repair.

When $p - q \equiv 1 \pmod 4$, adding one new direction of either sign fuses the halves: \[ \Cl(p{+}1,q) \;\cong\; \Cl(p,q{+}1) \;\cong\; \Mat\_2\bigl(\Cl\^{}0(p,q)\bigr) \] matrices over the even part. The fused algebra is simple, with centre $\R$, so it has no split of its own.

\begin{theorem}[Rule 2: ideal fusion, sharp form]\label{c:transport.rule-2-ideal-fusion-sharp-form}
Some Clifford algebras split into two independent halves. Adding one new direction, of either sign, fuses the halves into a single algebra of two-by-two matrices whose entries come from the algebra's even part. The fused algebra has no split of its own, and the two choices of sign give the same algebra.
\emph{First proved in} T1, Thm 5.5 (\href{https://doi.org/10.5281/zenodo.21743863}{doi:10.5281/zenodo.21743863}).
\end{theorem}
\noindent\emph{Rests on:} \ref{c:transport.ascent-renders-grade-involution-inner}, \ref{c:transport.centre}, \ref{c:transport.classification-real-clifford-algebras}, \ref{c:transport.corner-even-subalgebra}, \ref{c:transport.matrix-units-peirce-isomorphism}, \ref{c:std-algebra.morita-amplification}.

Is a single direction really the least one can add?

Yes: adding nothing leaves the algebra split, and adding a null direction only thickens it, by an ideal $J$ with $J^2 = 0$. One direction of nonzero square is the smallest repair that joins the halves, and both signs give the same algebra.

\begin{theorem}[One generator is minimal; the minimal completions classified]\label{c:transport.one-generator-minimal-minimal-completions-classified}
When a Clifford algebra splits into two halves, one can try to repair it by adding directions. Adding nothing leaves it split; adding one direction of zero length only thickens the algebra without joining the halves; adding one direction of nonzero length joins them; adding more is larger than needed. The two possible single directions, of positive or negative square, give the same algebra, so the repair is unique and cannot tell which sign was used.
\emph{First proved in} T1, Thm 5.7 (\href{https://doi.org/10.5281/zenodo.21743863}{doi:10.5281/zenodo.21743863}).
\end{theorem}
\noindent\emph{Rests on:} \ref{c:transport.decomposition-commutation}, \ref{c:transport.existence-classification-type-c}, \ref{c:transport.rule-2-ideal-fusion-sharp-form}, \ref{c:std-algebra.blade-basis}, \ref{c:std-algebra.sylvester-inertia}.

Repairs can be repeated, so it matters how often each is needed.

Every algebra sits on a clock, its residue $d \bmod 8$: Rule 1 never moves the hand, and Rule 2 is used only at residues $1$ and $5$ and lands in $\{0, 2\}$ or $\{4, 6\}$. So fusion happens at most once, and the algebra of spacetime starts at a residue that never needs it.

\begin{proposition}[Fusion occurs at most once]\label{c:transport.fusion-occurs-most-once}
Each Clifford algebra sits at one of eight positions on a clock. One repair rule never moves the hand; the other moves it by one step, is only used at the two positions where the algebra is split, and always lands where no split occurs. So the second rule can be used at most once, after which nothing structural changes. The algebra of spacetime starts at a position that never needs the second rule.
\emph{First proved in} T1, Prop 7.4 (\href{https://doi.org/10.5281/zenodo.21743863}{doi:10.5281/zenodo.21743863}).
\end{proposition}
\noindent\emph{Rests on:} \ref{c:transport.classification-real-clifford-algebras}, \ref{c:transport.existence-classification-type-c}, \ref{c:transport.rule-1-morita-trivial}.

Everything now fits into one statement.

Division in a nondegenerate real Clifford algebra fails in exactly three ways, and each carries a forced minimal repair: complete the space, complete the algebra, or do nothing. For a null vector the element itself cannot be repaired; the repair acts on the space of points, with increment $(+1,+1)$.

\begin{principle}[Transport principle]\label{c:transport.transport-principle}
Division can go wrong in a Clifford algebra in exactly three ways. For a light-like direction the element itself cannot be fixed, but adding one extra pair of dimensions repairs the space of points; for an algebra split into two halves, adding one dimension fuses them; for an ordinary projection, nothing needs fixing. Each repair is the smallest possible, and the list is complete.
\emph{First proved in} T1, Principle 8.2 (\href{https://doi.org/10.5281/zenodo.21743863}{doi:10.5281/zenodo.21743863}).
\end{principle}
\noindent\emph{Rests on:} \ref{c:transport.cohomological-form-no-go}, \ref{c:transport.embedding-independent-minimality-fusion}, \ref{c:transport.exhaustiveness-every-zero-divisor-annihilated-idempotent}, \ref{c:transport.minimality}, \ref{c:transport.no-fourth-rule}, \ref{c:transport.no-go-level-elements}, \ref{c:transport.one-generator-minimal-minimal-completions-classified}, \ref{c:transport.representation-independent-minimality-completion}, \ref{c:transport.rule-2-ideal-fusion-sharp-form}, \ref{c:transport.totality}, \ref{c:transport.type-p-self-repairing}.

\section{Results used from elsewhere}
The proofs above use the following results of the series, stated briefly; each is proved in the paper named.

\begin{fact}[The blade basis persists under degeneracy]\label{c:conformal.blade-basis-persists-under-degeneracy}
For any symmetric bilinear form on an $m$-dimensional real space, degenerate or not, the $2^m$ blades of any basis form a linear basis of $\Cl(W,B)$.
\emph{Source:} T2, Fact 2.1 (\href{https://doi.org/10.5281/zenodo.22714886}{doi:10.5281/zenodo.22714886}).
\end{fact}

\begin{fact}[The Witt extension theorem]\label{c:conformal.witt-extension-theorem}
Over a field of characteristic not two, an isometry between arbitrary subspaces of two isometric nondegenerate quadratic spaces extends to an isometry of the whole spaces.
\emph{Source:} T. Y. Lam, Introduction to Quadratic Forms over Fields, Graduate Studies in Mathematics 67, American Mathematical Society (2005)., Ch.~I, Thm.~4.7.
\end{fact}

\begin{fact}[The blade basis, degenerate forms included]\label{c:std-algebra.blade-basis}
For any basis of $V$ the $2^{n}$ increasing products $e_I$ form a basis of $\Cl(V, Q)$, degenerate forms included; so $V \hookrightarrow \Cl(V, Q)$.
\emph{Source:} the algebra core of the Transport Lab, proved there with an exact check; standard, see C. Chevalley, The Algebraic Theory of Spinors, Columbia University Press (1954), Ch. II (the Clifford algebra of a quadratic space).
\end{fact}

\begin{fact}[Isometries extend to algebra homomorphisms]\label{c:std-algebra.functoriality}
An isometry $f$ extends uniquely to an algebra homomorphism $\Cl(f)$, injective when $f$ is; $\mathrm{O}(V, Q)$ acts by automorphisms.
\emph{Source:} the algebra core of the Transport Lab, proved there with an exact check; standard, see P. Lounesto, Clifford Algebras and Spinors, 2nd ed., London Mathematical Society Lecture Note Series 286, Cambridge University Press (2001), Ch. 14 (definitions of the Clifford algebra).
\end{fact}

\begin{fact}[The three involutions on each grade]\label{c:std-algebra.involution-properties}
On $\Cl^{k}$: $\hat{x} = (-1)^{k}x$, $\tilde{x} = (-1)^{k(k-1)/2}x$, $\bar{x} = (-1)^{k(k+1)/2}x$; the first is an automorphism, the others anti-automorphisms.
\emph{Source:} the algebra core of the Transport Lab, proved there with an exact check; standard, see P. Lounesto, Clifford Algebras and Spinors, 2nd ed., London Mathematical Society Lecture Note Series 286, Cambridge University Press (2001), Ch. 14 (definitions of the Clifford algebra).
\end{fact}

\begin{fact}[An algebra and its matrix amplification]\label{c:std-algebra.morita-amplification}
$\Zc(\Mat_m(A)) = \Zc(A)$, the ideals of $\Mat_m(A)$ are $\Mat_m(I)$, and $A$ and $\Mat_m(A)$ have equivalent module categories.
\emph{Source:} the algebra core of the Transport Lab, proved there with an exact check; standard, see T. Y. Lam, Lectures on Modules and Rings, Graduate Texts in Mathematics 189, Springer, New York (1999), matrix rings, idempotents and Morita theory.
\end{fact}

\begin{fact}[Reflection in a non-null vector]\label{c:std-algebra.reflection}
For $Q(u) \ne 0$, $x \mapsto -uxu^{-1} = x - 2\frac{B(x,u)}{Q(u)}u$ is the reflection in $u^{\perp}$: an involutive isometry of determinant $-1$ fixing $\operatorname{rad} V$.
\emph{Source:} the algebra core of the Transport Lab, proved there with an exact check; standard, see C. Doran and A. Lasenby, Geometric Algebra for Physicists, Cambridge University Press (2003), Ch. 4 (foundations of geometric algebra: rotors, versors).
\end{fact}

\begin{fact}[Orthogonal bases and the law of inertia]\label{c:std-algebra.sylvester-inertia}
Every real quadratic space has an orthogonal basis squaring to $1, -1, 0$, and the counts $(p, q, r)$ are invariants that classify it up to isometry.
\emph{Source:} T. Y. Lam, Introduction to Quadratic Forms over Fields, AMS (2005)., the diagonalisation of forms and Sylvester's law of inertia (real case).
\end{fact}

\begin{fact}[Hyperbolic pairs and the Witt decomposition]\label{c:std-algebra.witt-decomposition}
A null vector in a nondegenerate plane has a unique hyperbolic partner, and $V = H_1 \perp \dots \perp H_{\min(p,q)} \perp A \perp \operatorname{rad} V$ with $A$ definite; maximal totally null subspaces have dimension $\min(p,q) + r$.
\emph{Source:} the algebra core of the Transport Lab, proved there with an exact check; standard, see T. Y. Lam, Introduction to Quadratic Forms over Fields, Graduate Studies in Mathematics 67, American Mathematical Society (2005), Ch. I (hyperbolic planes, the Witt decomposition, Witt's theorems).
\end{fact}

\begin{fact}[The ascent renders the grade involution inner]\label{c:transport.ascent-renders-grade-involution-inner}
In the one-generator ascent, conjugation by the adjoined generator $\nu$ is the grade involution on $A$, exchanging the projectors $\varepsilon_\pm$, so either one generates the whole algebra.
\emph{Source:} T1, Prop 5.2 (\href{https://doi.org/10.5281/zenodo.21743863}{doi:10.5281/zenodo.21743863}).
\end{fact}

\begin{fact}[The centre]\label{c:transport.centre}
For $n = p + q \geq 1$ the centre of $\Cl(p,q)$ is $\R \cdot 1$ when $n$ is even and $\R \cdot 1 \oplus \R \cdot \omega$ when $n$ is odd, with $\omega$ the pseudoscalar.
\emph{Source:} T1, Lem 2.6 (\href{https://doi.org/10.5281/zenodo.21743863}{doi:10.5281/zenodo.21743863}).
\end{fact}

\begin{fact}[Classification of the real Clifford algebras]\label{c:transport.classification-real-clifford-algebras}
The real Clifford algebra $\Cl(p,q)$ is a matrix algebra over $\R$, $\C$ or $\HH$, or a sum of two such, according to $p - q \bmod 8$, and it is simple except when $p - q \equiv 1 \pmod 4$.
\emph{Source:} M. F. Atiyah, R. Bott and A. Shapiro, Clifford modules (1964); I. R. Porteous, Clifford Algebras and the Classical Groups (1995); P. Lounesto, Clifford Algebras and Spinors (2001).
\end{fact}

\begin{fact}[Cohomological form of the no-go]\label{c:transport.cohomological-form-no-go}
The obstruction to an invariant pseudo-inverse of a null vector is the tautological class $\theta \in \mathrm{Hom}(N_v, U)$, which vanishes exactly when $U = 0$ and persists unchanged through every ascent and every larger symmetry group containing $N_v$.
\emph{Source:} T1, Thm 4.8 (\href{https://doi.org/10.5281/zenodo.21743863}{doi:10.5281/zenodo.21743863}).
\end{fact}

\begin{fact}[The corner is the even subalgebra]\label{c:transport.corner-even-subalgebra}
For $p - q \equiv 1 \pmod 4$, $a \mapsto \varepsilon_+ a$ is a unital isomorphism $\Cl^{0}(p,q) \to \varepsilon_+ A$, and $\varepsilon_+ B_\gamma \varepsilon_+ = \varepsilon_+ A$.
\emph{Source:} T1, Lem 5.3 (\href{https://doi.org/10.5281/zenodo.21743863}{doi:10.5281/zenodo.21743863}).
\end{fact}

\begin{fact}[Decomposition and commutation]\label{c:transport.decomposition-commutation}
The one-generator ascent splits as $B_\gamma = A \oplus A\nu$ with $\nu a = \hat a\, \nu$, and $\nu^{-1} = \gamma\nu$ when $\gamma = \pm 1$.
\emph{Source:} T1, Lem 5.1 (\href{https://doi.org/10.5281/zenodo.21743863}{doi:10.5281/zenodo.21743863}).
\end{fact}

\begin{fact}[Density]\label{c:transport.density}
The lifted invertible points $[X(x)]$ with $Q(x) \neq 0$ are dense in the quadric $\mathcal{Q}$.
\emph{Source:} T1, Lem 4.9 (\href{https://doi.org/10.5281/zenodo.21743863}{doi:10.5281/zenodo.21743863}).
\end{fact}

\begin{fact}[Embedding-independent minimality of the fusion]\label{c:transport.embedding-independent-minimality-fusion}
For $p - q \equiv 1 \pmod 4$, every simple real algebra containing $\Cl(p,q)$ unitally has dimension at least $2^{n+1}$, with equality only for $\Mat_2(\Cl^0(p,q))$.
\emph{Source:} T1, Thm 5.11 (\href{https://doi.org/10.5281/zenodo.21743863}{doi:10.5281/zenodo.21743863}).
\end{fact}

\begin{fact}[Exhaustiveness: every zero divisor is annihilated by an idempotent]\label{c:transport.exhaustiveness-every-zero-divisor-annihilated-idempotent}
Every nonzero zero divisor of a nondegenerate real Clifford algebra is annihilated on each side by a nontrivial idempotent, so the N, P, C taxonomy of obstructions is exhaustive.
\emph{Source:} T1, Thm 3.9 (\href{https://doi.org/10.5281/zenodo.21743863}{doi:10.5281/zenodo.21743863}).
\end{fact}

\begin{fact}[Existence and classification of type C]\label{c:transport.existence-classification-type-c}
$\Cl(p,q)$, $n \geq 1$, has a nontrivial central idempotent exactly when $p - q \equiv 1 \pmod 4$, and then the only ones are $\varepsilon_\pm = \tfrac12(1 \pm \omega)$, which split the algebra into two simple ideals and are exchanged by the grade involution but by no inner automorphism.
\emph{Source:} T1, Thm 3.6 (\href{https://doi.org/10.5281/zenodo.21743863}{doi:10.5281/zenodo.21743863}).
\end{fact}

\begin{fact}[Matrix units, and the Peirce isomorphism]\label{c:transport.matrix-units-peirce-isomorphism}
The projectors $\varepsilon_\pm$ and the adjoined generator $\nu$ form a full system of $2 \times 2$ matrix units, so $B_\gamma \cong \Mat_2(E_{11} B_\gamma E_{11})$.
\emph{Source:} T1, Lem 5.4 (\href{https://doi.org/10.5281/zenodo.21743863}{doi:10.5281/zenodo.21743863}).
\end{fact}

\begin{fact}[No fourth rule]\label{c:transport.no-fourth-rule}
Every zero divisor of a nondegenerate real Clifford algebra is of type N, C or P, handled respectively by a hyperbolic-plane ascent, a one-generator ascent, or no ascent, so there is no fourth rule.
\emph{Source:} T1, Cor 6.7 (\href{https://doi.org/10.5281/zenodo.21743863}{doi:10.5281/zenodo.21743863}).
\end{fact}

\begin{fact}[No-go at the level of elements]\label{c:transport.no-go-level-elements}
In mixed signature with $n \geq 3$, no assignment of von Neumann pseudo-inverses to null vectors, valued in the Clifford algebra of any extension, is equivariant even under $\Spin^0(p,q)$.
\emph{Source:} T1, Thm 4.2 (\href{https://doi.org/10.5281/zenodo.21743863}{doi:10.5281/zenodo.21743863}).
\end{fact}

\begin{fact}[Representation-independent minimality of the completion]\label{c:transport.representation-independent-minimality-completion}
Any nondegenerate quadratic space whose null cone carries a normalised conformal chart of $(\mathbb{V}, Q)$ has $p'' \geq p + 1$ and $q'' \geq q + 1$, so the minimality of the completion does not depend on how $\Cl(p,q)$ is embedded.
\emph{Source:} T1, Thm 4.19 (\href{https://doi.org/10.5281/zenodo.21743863}{doi:10.5281/zenodo.21743863}).
\end{fact}

\begin{fact}[Type P is self-repairing]\label{c:transport.type-p-self-repairing}
In any unital ring the identity assignment $\sigma(e) = e$ on idempotents satisfies both von Neumann identities and is equivariant under every automorphism of every extension preserving the ring, so type P obstructions are repaired at signature increment $(0,0)$.
\emph{Source:} T1, Prop 6.1 (\href{https://doi.org/10.5281/zenodo.21743863}{doi:10.5281/zenodo.21743863}).
\end{fact}

\section*{Conclusion}
The completion turns each of the three failures of division into a forced smallest repair.

A calculation done in the enlarged setting can be read back everywhere except at the points at infinity. This trail rests on no postulate: each step is proved here or cited from proved or classical results, such as the classification of real Clifford algebras and Witt's extension theorem.

\end{document}